\clearpage
\fancyhead[L]{\small HERSCHEL AND SOUTH. OBSERVATIONS OF DOUBLE STARS.}
\hypertarget{native-647-lower}{}
\begin{center}\small Göttingische gelehrte Anzeigen. Stück 201. Seite 2001..2004. 1826 December 18.\\[3pt]\rule{.72\linewidth}{.3pt}\end{center}

\emph{Observations of the apparent distances and positions of \textup{380} double and triple stars made in the years \textup{1821}, \textup{1822} and \textup{1823}, and compared with those of other astronomers.} London 1825. By \textsc{J. F. W. Herschel} and \textsc{James South}. 412 Seiten in 4., nebst 4 Kupfertafeln.

\emph{Observations of the apparent distances and positions of \textup{458} double and triple stars made in the years \textup{1823}, \textup{1824} and \textup{1825}.} Ebendaselbst 1826. By \textsc{James South}. 391 und XVIII Seiten in 4.

Die Geschichte der Astronomie lässt uns oft bemerken, \raisebox{.9ex}{.}dass wichtige unerwartete Aufschlüsse aus Beobachtungen erhalten wurden, die lange vorher mit mühsamem Fleiss, ohne einen solchen Erfolg zu ahnen, angestellt und aufgezeichnet waren. Die Doppelsterne sind ein merkwürdiges Beispiel dieser Art. Vor \textsc{Herschel} kannte man zwar schon eine, obwohl vergleichungsweise nur kleine, Anzahl von Doppelsternen, doch ohne sie einer besondern Aufmerksamkeit zu würdigen, indem man von der \SourceLeaf{648}Meinung ausging, dass das Phänomen eines Doppelsterns nur das Spiel eines Zufalls sei, der zwei Fixsterne mit unserm Sonnensysteme beinahe in einerlei Richtung gestellt habe, wobei die wirklichen Entfernungen im allgemeinen sehr ungleich sein möchten. \textsc{Herschel} machte zuerst auf die Vortheile aufmerksam, die die fortgesetzte Beobachtung eines Doppelsterns für die Bestimmung der jährlichen Parallaxe gewähren könnte, eine Methode, deren Brauchbarkeit offenbar ganz von jener Voraussetzung abhängig ist. \textsc{Herschel} zeichnete bei seiner Durchmusterung des gestirnten Himmels alle ihm vorkommenden Doppelsterne auf, beobachtete die gegenseitigen Stellungen mit vieler Genauigkeit, und brachte so ein Verzeichniss von mehr als 700 Doppelsternen und vielfachen Sternen zusammen, unter denen freilich viele sind, denen man wegen der beträchtlichen Entfernung der einzelnen Sterne von einander den Namen Doppelsterne nur uneigentlich beilegen kann, obwohl die Natur der Sache, in Rücksicht auf diese Entfernung, scharfe Grenzen zu ziehen nicht verstattet.

Zwanzig Jahre später unterwarf \textsc{Herschel} die Doppelsterne einer neuen Revision und bemerkte bei vielen derselben in der gegenseitigen Stellung der einzelnen Sterne solche Veränderungen, welche zu der Ueberzeugung führten, dass die frühere Ansicht einer zufälligen Bildung der Doppelsterne nicht allgemein zulässig, und dass wenigstens viele derselben \emph{wirklich} nahe zusammenstehende Sternpaare sein möchten, die nach bestimmten in den Naturkräften gegründeten Gesetzen sich um einander bewegen.

Diese höchst wichtige Folgerung erhielt durch die so sehr grosse Anzahl von Doppelsternen, welche am Himmel bemerkt werden, eine sehr verstärkte Wahrscheinlichkeit. Man kann aber dreist behaupten, dass es eigentlich alles dessen kaum bedurft hätte, um die frühere Voraussetzung in ihrer Blösse zu zeigen. In der That, wenn wir z. B. nur den einen Doppelstern, Castor, in Erwägung nehmen, welcher aus zwei Sternen besteht, die an Helligkeit nicht sehr viel verschieden, beide etwa zur dritten Grösse gehörig, ungefähr fünf Secunden von einander abstehen, so finden wir die Wahrscheinlichkeit der Bildung eines solchen Doppelsterns unter Voraussetzung einer ganz zufälligen Vertheilung sämmtlicher Fixsterne bis zur dritten Grösse auf der Himmelskugel so ausserordentlich klein, dass wir im gewöhnlichen Leben, bei Ereignissen, wo etwas ähnliches Statt findet, es für ungereimt halten, an einer \emph{nicht} zufälligen Entstehung zu zweifeln. Unmöglich ist es freilich nicht, dass unter der grossen Menge der Doppelsterne vielleicht einige von zufälliger Entstehung sein mögen; aber man muss nothwendig annehmen, dass $\cdot$die meisten der näher zusammenstehenden wirklich zusammengehörende Systeme bilden, deren gegenseitige Bewegungen ein neues Feld von interessanten Beobachtungen eröffnet haben. Diese Bewegungen sind aber so langsam, dass bei den meisten selbst in einem Menschenalter die Verrückung nur gering ist: doch haben sich auch einige Doppelsterne gefunden, wo dieselbe verhältnissmässig schon in sehr kurzen Fristen sichtbar wird, und auf alle Fälle ist klar, dass es um den künftigen Jahrhunderten den Stoff zu höchst erweiterter Kenntniss des Weltgebäudes vorzubereiten, von grosser Wichtigkeit ist, die Doppelsterne von Zeit zu Zeit einer neuen Revision zu unterwerfen.

Dies haben nun die Herren \textsc{Herschel} d. J., Sohn des berühmten Beobachters, und \textsc{South} in den vorliegenden Werken auf eine musterhafte Art gethan. Das erstere Werk enthält diejenigen Beobachtungen, welche diese beiden Astronomen gemeinschaftlich, mit einem fünffussigen und einem siebenfussigen Aequatoreal angestellt haben; das zweite diejenigen, welche der letztere Astronom allein, theils mit beiden Instrumenten in England, theils mit dem siebenfussigen Aequatoreal in Frankreich, in einer zu Passy interimistisch errichteten Sternwarte, gemacht hat. Wir haben hier also einen reichen Schatz von Beobachtungen über 838 Doppelsterne aus den Jahren 1821..1825, und zugleich ihre Zusammenstellung mit den frühern Beobachtungen des ältern \textsc{Herschel} zum Theil aus dem handschriftlichen Nachlass berichtigt und ergänzt, und den nicht minder schätzbaren von \textsc{Struve}. Bei den meisten Doppelsternen sind die bisher beobachteten Veränderungen in den Stellungen zwar nur klein, wenn gleich \SourceLeaf{649}entschieden: die Umlaufszeit einer Sonne um die andere wird im Allgemeinen nach Jahrhunderten gemessen. Inzwischen finden sich bei manchen schon so bedeutende Aenderungen, dass wir dadurch eine genäherte Vorstellung von der Grösse der Periode erhalten; bei dem merkwürdigen Sterne 61 Schwan z. B. hat die Veränderung des Richtungswinkels vom J. 1753 bis 1825 schon 53 Grad betragen; bei Castor von 1759 bis ebendahin 63 Grad; bei $\sigma$ Krone, seit 1781, 90 Grad; bei $\xi$ im grossen Bär in derselben Zwischenzeit 259 Grad; bei 70 Ophiuchus seit 1779 sogar 302 Grad. Der beschränkte Raum verbietet uns, noch manches andere Merkwürdige auszuheben, von Sternen, die früher doppelt erschienen und gegenwärtig einander vollkommen decken; von dreifachen Sternen, die Ein System bilden und sich gegenseitig bewegen; von Sternen, die zwölf Minuten von einander abstehend doch sehr wahrscheinlich nur Ein System bilden u. dergl. [S. Neue Aussicht zur Erweiterung u. s. f. oben S. 181. \textsc{S\raisebox{.9ex}{.}chering}.]
